% Options for packages loaded elsewhere
% Options for packages loaded elsewhere
\PassOptionsToPackage{unicode}{hyperref}
\PassOptionsToPackage{hyphens}{url}
%
\documentclass[
  ignorenonframetext,
  aspectratio=169,
  spanish,
]{beamer}
\newif\ifbibliography
\usepackage{pgfpages}
\setbeamertemplate{caption}[numbered]
\setbeamertemplate{caption label separator}{: }
\setbeamercolor{caption name}{fg=normal text.fg}
\beamertemplatenavigationsymbolsempty
% remove section numbering
\setbeamertemplate{part page}{
  \centering
  \begin{beamercolorbox}[sep=16pt,center]{part title}
    \usebeamerfont{part title}\insertpart\par
  \end{beamercolorbox}
}
\setbeamertemplate{section page}{
  \centering
  \begin{beamercolorbox}[sep=12pt,center]{section title}
    \usebeamerfont{section title}\insertsection\par
  \end{beamercolorbox}
}
\setbeamertemplate{subsection page}{
  \centering
  \begin{beamercolorbox}[sep=8pt,center]{subsection title}
    \usebeamerfont{subsection title}\insertsubsection\par
  \end{beamercolorbox}
}
% Prevent slide breaks in the middle of a paragraph
\widowpenalties 1 10000
\raggedbottom
\AtBeginPart{
  \frame{\partpage}
}
\AtBeginSection{
  \ifbibliography
  \else
    \frame{\sectionpage}
  \fi
}
\AtBeginSubsection{
  \frame{\subsectionpage}
}
\usepackage{iftex}
\ifPDFTeX
  \usepackage[T1]{fontenc}
  \usepackage[utf8]{inputenc}
  \usepackage{textcomp} % provide euro and other symbols
\else % if luatex or xetex
  \usepackage{unicode-math} % this also loads fontspec
  \defaultfontfeatures{Scale=MatchLowercase}
  \defaultfontfeatures[\rmfamily]{Ligatures=TeX,Scale=1}
\fi
\usepackage{lmodern}

\ifPDFTeX\else
  % xetex/luatex font selection
\fi
% Use upquote if available, for straight quotes in verbatim environments
\IfFileExists{upquote.sty}{\usepackage{upquote}}{}
\IfFileExists{microtype.sty}{% use microtype if available
  \usepackage[]{microtype}
  \UseMicrotypeSet[protrusion]{basicmath} % disable protrusion for tt fonts
}{}
\makeatletter
\@ifundefined{KOMAClassName}{% if non-KOMA class
  \IfFileExists{parskip.sty}{%
    \usepackage{parskip}
  }{% else
    \setlength{\parindent}{0pt}
    \setlength{\parskip}{6pt plus 2pt minus 1pt}}
}{% if KOMA class
  \KOMAoptions{parskip=half}}
\makeatother

\usepackage{color}
\usepackage{fancyvrb}
\newcommand{\VerbBar}{|}
\newcommand{\VERB}{\Verb[commandchars=\\\{\}]}
\DefineVerbatimEnvironment{Highlighting}{Verbatim}{commandchars=\\\{\}}
% Add ',fontsize=\small' for more characters per line
\usepackage{framed}
\definecolor{shadecolor}{RGB}{241,243,245}
\newenvironment{Shaded}{\begin{snugshade}}{\end{snugshade}}
\newcommand{\AlertTok}[1]{\textcolor[rgb]{0.68,0.00,0.00}{#1}}
\newcommand{\AnnotationTok}[1]{\textcolor[rgb]{0.37,0.37,0.37}{#1}}
\newcommand{\AttributeTok}[1]{\textcolor[rgb]{0.40,0.45,0.13}{#1}}
\newcommand{\BaseNTok}[1]{\textcolor[rgb]{0.68,0.00,0.00}{#1}}
\newcommand{\BuiltInTok}[1]{\textcolor[rgb]{0.00,0.23,0.31}{#1}}
\newcommand{\CharTok}[1]{\textcolor[rgb]{0.13,0.47,0.30}{#1}}
\newcommand{\CommentTok}[1]{\textcolor[rgb]{0.37,0.37,0.37}{#1}}
\newcommand{\CommentVarTok}[1]{\textcolor[rgb]{0.37,0.37,0.37}{\textit{#1}}}
\newcommand{\ConstantTok}[1]{\textcolor[rgb]{0.56,0.35,0.01}{#1}}
\newcommand{\ControlFlowTok}[1]{\textcolor[rgb]{0.00,0.23,0.31}{\textbf{#1}}}
\newcommand{\DataTypeTok}[1]{\textcolor[rgb]{0.68,0.00,0.00}{#1}}
\newcommand{\DecValTok}[1]{\textcolor[rgb]{0.68,0.00,0.00}{#1}}
\newcommand{\DocumentationTok}[1]{\textcolor[rgb]{0.37,0.37,0.37}{\textit{#1}}}
\newcommand{\ErrorTok}[1]{\textcolor[rgb]{0.68,0.00,0.00}{#1}}
\newcommand{\ExtensionTok}[1]{\textcolor[rgb]{0.00,0.23,0.31}{#1}}
\newcommand{\FloatTok}[1]{\textcolor[rgb]{0.68,0.00,0.00}{#1}}
\newcommand{\FunctionTok}[1]{\textcolor[rgb]{0.28,0.35,0.67}{#1}}
\newcommand{\ImportTok}[1]{\textcolor[rgb]{0.00,0.46,0.62}{#1}}
\newcommand{\InformationTok}[1]{\textcolor[rgb]{0.37,0.37,0.37}{#1}}
\newcommand{\KeywordTok}[1]{\textcolor[rgb]{0.00,0.23,0.31}{\textbf{#1}}}
\newcommand{\NormalTok}[1]{\textcolor[rgb]{0.00,0.23,0.31}{#1}}
\newcommand{\OperatorTok}[1]{\textcolor[rgb]{0.37,0.37,0.37}{#1}}
\newcommand{\OtherTok}[1]{\textcolor[rgb]{0.00,0.23,0.31}{#1}}
\newcommand{\PreprocessorTok}[1]{\textcolor[rgb]{0.68,0.00,0.00}{#1}}
\newcommand{\RegionMarkerTok}[1]{\textcolor[rgb]{0.00,0.23,0.31}{#1}}
\newcommand{\SpecialCharTok}[1]{\textcolor[rgb]{0.37,0.37,0.37}{#1}}
\newcommand{\SpecialStringTok}[1]{\textcolor[rgb]{0.13,0.47,0.30}{#1}}
\newcommand{\StringTok}[1]{\textcolor[rgb]{0.13,0.47,0.30}{#1}}
\newcommand{\VariableTok}[1]{\textcolor[rgb]{0.07,0.07,0.07}{#1}}
\newcommand{\VerbatimStringTok}[1]{\textcolor[rgb]{0.13,0.47,0.30}{#1}}
\newcommand{\WarningTok}[1]{\textcolor[rgb]{0.37,0.37,0.37}{\textit{#1}}}

\usepackage{longtable,booktabs,array}
\newcounter{none} % for unnumbered tables
\usepackage{calc} % for calculating minipage widths
\usepackage{caption}
% Make caption package work with longtable
\makeatletter
\def\fnum@table{\tablename~\thetable}
\makeatother
\usepackage{graphicx}
\makeatletter
\newsavebox\pandoc@box
\newcommand*\pandocbounded[1]{% scales image to fit in text height/width
  \sbox\pandoc@box{#1}%
  \Gscale@div\@tempa{\textheight}{\dimexpr\ht\pandoc@box+\dp\pandoc@box\relax}%
  \Gscale@div\@tempb{\linewidth}{\wd\pandoc@box}%
  \ifdim\@tempb\p@<\@tempa\p@\let\@tempa\@tempb\fi% select the smaller of both
  \ifdim\@tempa\p@<\p@\scalebox{\@tempa}{\usebox\pandoc@box}%
  \else\usebox{\pandoc@box}%
  \fi%
}
% Set default figure placement to htbp
\def\fps@figure{htbp}
\makeatother



\ifLuaTeX
\usepackage[bidi=basic,shorthands=off]{babel}
\else
\usepackage[bidi=default,shorthands=off]{babel}
\fi
\ifLuaTeX
  \usepackage{selnolig} % disable illegal ligatures
\fi


\setlength{\emergencystretch}{3em} % prevent overfull lines

\providecommand{\tightlist}{%
  \setlength{\itemsep}{0pt}\setlength{\parskip}{0pt}}



 


\usepackage{booktabs}
\usepackage{amsmath}
\usepackage{amssymb}
\usepackage{mathtools}
\usepackage{xcolor}
\definecolor{crossblue}{RGB}{31,78,121}
\definecolor{crossgreen}{RGB}{0,120,110}
\setbeamercolor{structure}{fg=crossblue}
\setbeamercolor{alerted text}{fg=crossgreen}
\makeatletter
\@ifpackageloaded{caption}{}{\usepackage{caption}}
\AtBeginDocument{%
\ifdefined\contentsname
  \renewcommand*\contentsname{Tabla de contenidos}
\else
  \newcommand\contentsname{Tabla de contenidos}
\fi
\ifdefined\listfigurename
  \renewcommand*\listfigurename{Listado de Figuras}
\else
  \newcommand\listfigurename{Listado de Figuras}
\fi
\ifdefined\listtablename
  \renewcommand*\listtablename{Listado de Tablas}
\else
  \newcommand\listtablename{Listado de Tablas}
\fi
\ifdefined\figurename
  \renewcommand*\figurename{Figura}
\else
  \newcommand\figurename{Figura}
\fi
\ifdefined\tablename
  \renewcommand*\tablename{Tabla}
\else
  \newcommand\tablename{Tabla}
\fi
}
\@ifpackageloaded{float}{}{\usepackage{float}}
\floatstyle{ruled}
\@ifundefined{c@chapter}{\newfloat{codelisting}{h}{lop}}{\newfloat{codelisting}{h}{lop}[chapter]}
\floatname{codelisting}{Listado}
\newcommand*\listoflistings{\listof{codelisting}{Listado de Listados}}
\makeatother
\makeatletter
\makeatother
\makeatletter
\@ifpackageloaded{caption}{}{\usepackage{caption}}
\@ifpackageloaded{subcaption}{}{\usepackage{subcaption}}
\makeatother

\usepackage{bookmark}
\IfFileExists{xurl.sty}{\usepackage{xurl}}{} % add URL line breaks if available
\urlstyle{same}
% fallback for those not using the hyperref driver hyperxmp:
\makeatletter
\@ifundefined{xmpquote}{\newcommand{\xmpquote}[1]{#1}}{}
\makeatother
\hypersetup{
  pdftitle={¿Qué es CrossCarry?},
  pdfauthor={Nelson Alirio Cruz},
  pdflang={es},
  hidelinks,
  pdfcreator={LaTeX via pandoc}}


\title{\texorpdfstring{¿Qué es CrossCarry?}{¿Qué es CrossCarry?}}
\subtitle{\texorpdfstring{Innovación estadística en R para el análisis
de diseños crossover en salud y producción
animal}{Innovación estadística en R para el análisis de diseños crossover en salud y producción animal}}
\author{Nelson Alirio Cruz}
\date{Invalid Date}
\institute{Facultad de Medicina Veterinaria y Zootecnia · UNAM}

\begin{document}
\frame{\titlepage}


\section{¿Qué es CrossCarry?}\label{quuxe9-es-crosscarry}

\begin{frame}{¿Qué es CrossCarry?}
Una propuesta estadística y computacional para analizar diseños
crossover con respuestas correlacionadas, medidas repetidas y efectos
carry-over.

\vfill

\centering

\textbf{Nelson Alirio Cruz}\\
Facultad de Medicina Veterinaria y Zootecnia · UNAM\\
27 de octubre de 2026
\end{frame}

\begin{frame}{1. ¿Por qué estudiar diseños crossover?}
\protect\phantomsection\label{por-quuxe9-estudiar-diseuxf1os-crossover}
Los diseños crossover permiten comparar tratamientos aplicados
sucesivamente a las mismas unidades experimentales. Cada persona, animal
o unidad recibe más de un tratamiento, por lo que su respuesta bajo una
intervención puede compararse con su propia respuesta bajo otra.

\begin{itemize}
\tightlist
\item
  Reducen la variabilidad debida a diferencias permanentes entre
  individuos.
\item
  Pueden alcanzar buena precisión con menos unidades que un diseño
  paralelo.
\item
  Son útiles cuando la respuesta puede medirse repetidamente y el estado
  de la unidad es relativamente estable.
\item
  Son frecuentes en ensayos clínicos, nutrición animal, fisiología,
  comportamiento y estudios de ambientes de trabajo.
\end{itemize}

\alert{La ventaja principal es también la fuente del problema:} como una
misma unidad recibe tratamientos sucesivos, las respuestas observadas en
distintos períodos no son independientes.
\end{frame}

\begin{frame}{2. Anatomía de un diseño crossover}
\protect\phantomsection\label{anatomuxeda-de-un-diseuxf1o-crossover}
Un diseño crossover contiene cuatro elementos básicos:

\begin{itemize}
\tightlist
\item
  \textbf{Tratamientos:} intervenciones que se desean comparar.
\item
  \textbf{Períodos:} intervalos durante los cuales se aplica cada
  tratamiento.
\item
  \textbf{Secuencias:} órdenes posibles de administración de los
  tratamientos.
\item
  \textbf{Unidades experimentales:} individuos o elementos asignados a
  una secuencia.
\end{itemize}

En el diseño AB/BA más sencillo:

\[
\begin{array}{c|cc}
\text{Secuencia} & \text{Período 1} & \text{Período 2}\\ \hline
AB & A & B\\
BA & B & A
\end{array}
\]

La comparación no es solamente entre \(A\) y \(B\). También deben
considerarse el período, el orden de administración y la posible
persistencia del tratamiento anterior.
\end{frame}

\begin{frame}{3. Diseño paralelo, longitudinal y crossover}
\protect\phantomsection\label{diseuxf1o-paralelo-longitudinal-y-crossover}
En un diseño \textbf{paralelo}, cada unidad recibe un único tratamiento.
La comparación principal se realiza entre grupos diferentes.

En un estudio \textbf{longitudinal}, cada unidad se observa varias
veces, generalmente bajo el mismo tratamiento.

En un diseño \textbf{crossover}, cada unidad:

\begin{enumerate}
\tightlist
\item
  recibe una secuencia de tratamientos;
\item
  aporta observaciones en varios períodos;
\item
  puede tener varias mediciones dentro de cada período;
\item
  puede conservar efectos de tratamientos previos.
\end{enumerate}

Por tanto, un crossover con medidas repetidas combina dos dependencias:

\[
\text{dependencia longitudinal}
+
\text{dependencia por secuencia de tratamientos}.
\]

Esta doble estructura explica por qué un análisis convencional de medias
o una regresión que suponga independencia puede resultar insuficiente.
\end{frame}

\begin{frame}{4. Ventajas y condiciones de aplicación}
\protect\phantomsection\label{ventajas-y-condiciones-de-aplicaciuxf3n}
Un crossover resulta especialmente atractivo cuando:

\begin{itemize}
\tightlist
\item
  el efecto del tratamiento es reversible o puede disminuir
  suficientemente;
\item
  la enfermedad o condición se mantiene relativamente estable;
\item
  la unidad puede recibir todas las intervenciones sin riesgos
  indebidos;
\item
  la duración total del estudio es aceptable;
\item
  existe interés en comparaciones dentro de la misma unidad.
\end{itemize}

No es ideal cuando:

\begin{itemize}
\tightlist
\item
  el tratamiento produce curación permanente o cambios irreversibles;
\item
  la condición evoluciona rápidamente;
\item
  existe mortalidad o abandono diferencial importante;
\item
  la exposición anterior modifica de manera permanente la respuesta
  futura;
\item
  el tratamiento solo puede administrarse una vez.
\end{itemize}

\alert{El crossover no es solo una forma eficiente de recolectar datos. Es una estrategia de diseño que exige pensar desde el comienzo en período, secuencia, dependencia y carry-over.}
\end{frame}

\begin{frame}{5. ¿Qué es un período de lavado?}
\protect\phantomsection\label{quuxe9-es-un-peruxedodo-de-lavado}
El \textbf{período de lavado} o \emph{washout} es el intervalo entre dos
tratamientos consecutivos durante el cual se espera que desaparezca el
efecto de la intervención anterior.

\[
A
\longrightarrow
\underbrace{\text{lavado}}_{\text{sin tratamiento o condición neutra}}
\longrightarrow
B
\]

Su duración debe estar fundamentada científicamente. En farmacología
puede relacionarse con la vida media y la eliminación del fármaco. En
nutrición o producción animal puede depender de la adaptación
metabólica, el tránsito digestivo o el tiempo necesario para estabilizar
la producción.

Un lavado adecuado busca que, al comenzar el nuevo período,

\[
\text{efecto residual del tratamiento previo}\approx 0.
\]

Sin embargo, la existencia formal de un lavado no garantiza que el
carry-over haya desaparecido.
\end{frame}

\begin{frame}{6. ¿Cuándo el lavado es imposible o insuficiente?}
\protect\phantomsection\label{cuuxe1ndo-el-lavado-es-imposible-o-insuficiente}
En muchas aplicaciones no es posible suspender completamente la
intervención:

\begin{itemize}
\tightlist
\item
  pacientes con hipertensión no pueden permanecer largos períodos sin
  tratamiento;
\item
  animales deben continuar alimentándose todos los días;
\item
  dietas y rutinas de manejo generan adaptación progresiva;
\item
  un placebo no necesariamente elimina el efecto acumulado de una
  intervención previa;
\item
  los ambientes de trabajo o alojamiento pueden producir adaptación
  conductual;
\item
  por razones éticas, no siempre puede retirarse agua, alimento o
  atención.
\end{itemize}

En estos casos, el análisis no debe utilizar el lavado como una
justificación automática para eliminar el carry-over del modelo.

\alert{CrossCarry no reemplaza un buen período de lavado.} Proporciona
herramientas para evaluar y modelar efectos residuales cuando no pueden
descartarse por diseño.
\end{frame}

\begin{frame}{7. ¿Qué es el efecto carry-over?}
\protect\phantomsection\label{quuxe9-es-el-efecto-carry-over}
El carry-over es la influencia persistente de un tratamiento
administrado previamente sobre una respuesta observada bajo un
tratamiento posterior.

Si una unidad recibe \(A\) en el período 1 y \(B\) en el período 2, la
respuesta del segundo período puede descomponerse conceptualmente como

\[
\text{respuesta bajo }B
=
\text{efecto directo de }B
+
\text{residuo de }A
+
\text{otros efectos}.
\]

Se habla de carry-over de primer orden cuando procede del período
inmediatamente anterior. Un efecto de segundo orden procede de dos
períodos atrás.

El carry-over puede afectar:

\begin{itemize}
\tightlist
\item
  el nivel promedio de la respuesta;
\item
  la forma temporal de la respuesta;
\item
  la variabilidad;
\item
  la estimación del tratamiento actual;
\item
  los errores estándar y las pruebas de hipótesis.
\end{itemize}
\end{frame}

\begin{frame}{8. Carry-over simple y complejo}
\protect\phantomsection\label{carry-over-simple-y-complejo}
En el carry-over simple, el efecto residual depende únicamente del
tratamiento previo:

\[
A\to B=A\to C=A\to D.
\]

En el carry-over complejo, depende de la transición completa:

\[
f_{A\to B}(t),\qquad
f_{A\to C}(t),\qquad
f_{A\to D}(t).
\]

Así, el residuo de \(A\) puede aumentar la respuesta bajo \(B\),
disminuirla bajo \(C\) y ser despreciable bajo \(D\).

Con \(G\) tratamientos pueden existir hasta

\[
C=G(G-1)
\]

transiciones ordenadas entre tratamientos distintos.

\alert{El carry-over complejo representa una interacción entre el tratamiento previo, el tratamiento actual y, cuando hay medidas repetidas, el tiempo intraperíodo.}
\end{frame}

\begin{frame}{9. Ejemplo AB/BA con una observación por período}
\protect\phantomsection\label{ejemplo-abba-con-una-observaciuxf3n-por-peruxedodo}
En un diseño AB/BA se observan dos transiciones:

\begin{itemize}
\tightlist
\item
  \(A\to B\) en el segundo período de la secuencia AB;
\item
  \(B\to A\) en el segundo período de la secuencia BA.
\end{itemize}

Con una sola observación por período es posible representar un
carry-over constante, pero no estimar de manera flexible una curva
temporal.

Si se desea ajustar

\[
f_{A\to B}(t)
=
\sum_{\ell=1}^{d}\theta_{A\to B,\ell}\phi_\ell(t),
\]

se requieren varios tiempos distintos. Una observación no puede
identificar una base con \(d>1\) coeficientes.

Este ejemplo separa dos ideas:

\begin{itemize}
\tightlist
\item
  la transición puede estar presente en el diseño;
\item
  la información puede ser insuficiente para estimar su evolución
  temporal.
\end{itemize}
\end{frame}

\begin{frame}{10. Ejemplo AB/BA con medidas repetidas}
\protect\phantomsection\label{ejemplo-abba-con-medidas-repetidas}
Supongamos ahora que cada unidad aporta \(L=10\) observaciones dentro de
cada período. Las transiciones \(A\to B\) y \(B\to A\) se observan a lo
largo de diez tiempos intraperíodo.

\[
\mathbf Y_{ij}
=
(Y_{ij1},\ldots,Y_{ij10})^\top.
\]

Esto permite estudiar:

\begin{itemize}
\tightlist
\item
  la tendencia común dentro del período;
\item
  la evolución del efecto residual de \(A\) sobre \(B\);
\item
  la evolución del efecto residual de \(B\) sobre \(A\);
\item
  la correlación entre mediciones cercanas;
\item
  la dependencia entre los dos períodos.
\end{itemize}

La ganancia de información no procede solamente de aumentar el número
total de filas, sino de observar la respuesta en diferentes momentos
relevantes del proceso.
\end{frame}

\begin{frame}{11. Diseño de Williams con cuatro tratamientos}
\protect\phantomsection\label{diseuxf1o-de-williams-con-cuatro-tratamientos}
Un diseño de Williams distribuye los tratamientos de forma equilibrada
entre períodos y secuencias. Un ejemplo de cuatro tratamientos es:

\[
\begin{array}{c|cccc}
\text{Secuencia} & P_1 & P_2 & P_3 & P_4\\ \hline
1 & B & A & D & C\\
2 & C & D & A & B\\
3 & D & B & C & A\\
4 & A & C & B & D
\end{array}
\]

En este arreglo aparecen las doce transiciones entre tratamientos
distintos. Por ejemplo:

\begin{itemize}
\tightlist
\item
  \(B\to A\), \(A\to D\) y \(D\to C\) en la primera secuencia;
\item
  \(C\to D\), \(D\to A\) y \(A\to B\) en la segunda.
\end{itemize}

Con medidas repetidas dentro de cada período, este diseño puede
proporcionar soporte para las doce funciones de carry-over complejo.
\end{frame}

\begin{frame}{12. Diseño incompleto y transiciones no observadas}
\protect\phantomsection\label{diseuxf1o-incompleto-y-transiciones-no-observadas}
Considérese un estudio con tres tratamientos y solo dos secuencias:

\[
A\to B\to C,
\qquad
B\to A\to C.
\]

Se observan las transiciones:

\[
A\to B,\quad B\to C,\quad B\to A,\quad A\to C.
\]

No aparecen \(C\to A\) ni \(C\to B\). Por tanto, sus efectos no pueden
estimarse, aunque existan muchos sujetos o muchas mediciones dentro de
los períodos observados.

\alert{Aumentar el tamaño muestral no crea una transición ausente.}

Antes del análisis debe construirse una tabla de transiciones y
comprobar:

\begin{itemize}
\tightlist
\item
  cuáles aparecen;
\item
  en qué períodos aparecen;
\item
  cuántas unidades las aportan;
\item
  cuántos tiempos intraperíodo las respaldan.
\end{itemize}
\end{frame}

\begin{frame}{13. Notación para medidas repetidas}
\protect\phantomsection\label{notaciuxf3n-para-medidas-repetidas}
Sea \(Y_{ijk}\) la respuesta de la unidad \(i\), en el período \(j\),
durante la medición intraperíodo \(k\).

\[
\mathbf Y_{ij}
=(Y_{ij1},\ldots,Y_{ijL})^\top,
\]

\[
\mathbf Y_i
=(\mathbf Y_{i1}^\top,\ldots,\mathbf Y_{iP}^\top)^\top.
\]

El índice \(t_{ijk}\) representa el tiempo de medición dentro del
período, no necesariamente el tiempo absoluto desde el inicio del
estudio.

En algunas aplicaciones farmacológicas es preferible redefinirlo como
tiempo desde la exposición previa. La formulación estadística se
conserva, pero cambia la interpretación científica de las funciones.

Además de la respuesta, el conjunto de datos debe contener:

\begin{itemize}
\tightlist
\item
  identificador de unidad;
\item
  secuencia;
\item
  período;
\item
  tratamiento actual;
\item
  tratamiento anterior o transición;
\item
  tiempo intraperíodo;
\item
  covariables relevantes.
\end{itemize}
\end{frame}

\begin{frame}[fragile]{14. Datos en formato largo}
\protect\phantomsection\label{datos-en-formato-largo}
CrossCarry trabaja naturalmente con una fila por observación. La
estructura mínima puede inspeccionarse así:

\begin{Shaded}
\begin{Highlighting}[]
\CommentTok{\# Una fila por unidad, período y tiempo}
\FunctionTok{head}\NormalTok{(datos[, }\FunctionTok{c}\NormalTok{(}
  \StringTok{"id"}\NormalTok{, }\StringTok{"secuencia"}\NormalTok{, }\StringTok{"periodo"}\NormalTok{,}
  \StringTok{"tratamiento"}\NormalTok{, }\StringTok{"tiempo"}\NormalTok{, }\StringTok{"respuesta"}
\NormalTok{)])}
\end{Highlighting}
\end{Shaded}

Es indispensable verificar el orden antes de construir las variables
carry-over:

\begin{Shaded}
\begin{Highlighting}[]
\NormalTok{datos }\OtherTok{\textless{}{-}}\NormalTok{ datos }\SpecialCharTok{|\textgreater{}}
\NormalTok{  dplyr}\SpecialCharTok{::}\FunctionTok{arrange}\NormalTok{(id, periodo, tiempo)}
\end{Highlighting}
\end{Shaded}

Controles recomendados:

\begin{Shaded}
\begin{Highlighting}[]
\FunctionTok{with}\NormalTok{(datos, }\FunctionTok{table}\NormalTok{(secuencia, periodo, tratamiento))}
\FunctionTok{with}\NormalTok{(datos, }\FunctionTok{table}\NormalTok{(id, periodo))}
\FunctionTok{summary}\NormalTok{(datos}\SpecialCharTok{$}\NormalTok{tiempo)}
\end{Highlighting}
\end{Shaded}

Un error de ordenamiento puede asignar como tratamiento previo una
observación equivocada y alterar toda la interpretación del modelo.
\end{frame}

\begin{frame}{15. Condiciones básicas de estimabilidad}
\protect\phantomsection\label{condiciones-buxe1sicas-de-estimabilidad}
Para separar tratamiento, período y carry-over deben cumplirse
condiciones estructurales:

\begin{enumerate}
\tightlist
\item
  cada tratamiento debe aparecer en más de una secuencia;
\item
  cada tratamiento debe aparecer en más de un período;
\item
  cada período debe contener más de un tratamiento;
\item
  cada transición incluida debe observarse en al menos una
  unidad-período;
\item
  las columnas asociadas con tratamiento, período y carry-over no deben
  ser linealmente dependientes.
\end{enumerate}

Estas condiciones se expresan mediante el rango de la matriz de diseño:

\[
\operatorname{rank}(\mathbf X)
=
\text{número de parámetros identificables}.
\]

La estimabilidad es una propiedad conjunta del diseño, la codificación
de los efectos y las funciones de base elegidas.
\end{frame}

\begin{frame}{16. Estimabilidad de las funciones}
\protect\phantomsection\label{estimabilidad-de-las-funciones}
Si un efecto funcional se representa con una base de dimensión \(d\),

\[
f_c(t)
=
\sum_{\ell=1}^{d}\theta_{c\ell}\phi_\ell(t),
\]

la matriz de evaluación debe satisfacer

\[
\operatorname{rank}(\boldsymbol\Phi_c)=d,
\qquad L\ge d.
\]

No basta con tener \(L\) observaciones si todas se concentran en pocos
tiempos o si la base genera columnas dependientes.

En la práctica:

\begin{itemize}
\tightlist
\item
  una base pequeña suele ser preferible con pocos tiempos;
\item
  la dimensión debe crecer de manera conservadora;
\item
  cada transición necesita replicación entre unidades, no solo muchos
  registros del mismo individuo;
\item
  la penalización mejora estabilidad, pero no corrige falta total de
  información.
\end{itemize}

\alert{Identificabilidad teórica y estimabilidad práctica no son sinónimos.}
\end{frame}

\begin{frame}{17. Primera etapa metodológica: respuestas de conteo}
\protect\phantomsection\label{primera-etapa-metodoluxf3gica-respuestas-de-conteo}
Uno de los desarrollos que fundamenta CrossCarry estudió respuestas
Poisson en diseños crossover. Para una media \(\mu_{ij}\),

\[
Y_{ijk}\sim\operatorname{Poisson}(\mu_{ij}),
\]

\[
\log(\mu_{ij})
=
\mu+\tau_{ij}+\pi_j+\theta_{ij}.
\]

El modelo separa:

\begin{itemize}
\tightlist
\item
  media global;
\item
  tratamiento;
\item
  período;
\item
  carry-over.
\end{itemize}

Para diseños pequeños, como AB/BA, algunos estimadores pueden expresarse
mediante combinaciones de logaritmos de medias de celda. El método delta
permite aproximar sus varianzas.

Este desarrollo mostró que los modelos normales convencionales pueden
producir estimaciones deficientes cuando la respuesta es un conteo y
existe carry-over.
\end{frame}

\begin{frame}{18. Sobredispersión y subdispersión}
\protect\phantomsection\label{sobredispersiuxf3n-y-subdispersiuxf3n}
La distribución Poisson supone igualdad entre media y varianza:

\[
\operatorname{Var}(Y_{ijk})=E(Y_{ijk})=\mu_{ijk}.
\]

En datos reales puede aparecer dispersión adicional:

\[
\operatorname{Var}(Y_{ijk})
=
\phi\mu_{ijk}.
\]

\begin{itemize}
\tightlist
\item
  \(\phi>1\): sobredispersión;
\item
  \(\phi<1\): subdispersión;
\item
  \(\phi=1\): comportamiento Poisson ideal.
\end{itemize}

Ignorar la dispersión afecta principalmente los errores estándar y las
pruebas de hipótesis. Por ello, la metodología incorporó un parámetro de
dispersión y posteriormente se extendió mediante GEE a diseños con más
períodos, covariables y estructuras de correlación.

\alert{La distribución de la respuesta y la dependencia entre observaciones son problemas diferentes y ambos deben modelarse.}
\end{frame}

\begin{frame}{19. Aplicación: hidratación y cognición en Mali}
\protect\phantomsection\label{aplicaciuxf3n-hidrataciuxf3n-y-cogniciuxf3n-en-mali}
En el estudio de hidratación escolar, 47 estudiantes recibieron primero
control y después hidratación, mientras 60 recibieron el orden inverso.
La respuesta fue el número de diferencias correctamente identificadas
entre dos imágenes.

Características del problema:

\begin{itemize}
\tightlist
\item
  respuesta de conteo;
\item
  diseño AB/BA;
\item
  ausencia de un lavado realista por razones éticas;
\item
  grado escolar como covariable;
\item
  observaciones correlacionadas dentro de cada estudiante.
\end{itemize}

El modelo GEE con grado escolar presentó mejor QIC que el análisis
loglineal más sencillo. Los resultados indicaron:

\begin{itemize}
\tightlist
\item
  efecto de período;
\item
  efecto de grado;
\item
  carry-over significativo;
\item
  ausencia de un efecto directo significativo de hidratación en el
  período de aplicación.
\end{itemize}

El ejemplo muestra que un tratamiento puede influir posteriormente
aunque su efecto directo inmediato no resulte significativo.
\end{frame}

\begin{frame}{20. ¿Por qué usar GEE?}
\protect\phantomsection\label{por-quuxe9-usar-gee}
Las ecuaciones de estimación generalizadas modelan la media marginal y
reconocen la correlación dentro de cada unidad.

La ecuación básica es

\[
\mathbf U(\boldsymbol\beta)
=
\sum_{i=1}^{n}
\mathbf D_i^\top
\mathbf V_i^{-1}
(\mathbf Y_i-\boldsymbol\mu_i)
=
\mathbf 0.
\]

Aquí:

\begin{itemize}
\tightlist
\item
  \(\boldsymbol\mu_i\) contiene las medias marginales;
\item
  \(\mathbf D_i=\partial\boldsymbol\mu_i/\partial\boldsymbol\beta^\top\);
\item
  \(\mathbf V_i\) es una matriz de covarianza de trabajo.
\end{itemize}

Ventajas:

\begin{itemize}
\tightlist
\item
  permite distintas familias y enlaces;
\item
  produce efectos marginales interpretables;
\item
  usa varianza robusta tipo sándwich;
\item
  no exige especificar una distribución completa para efectos
  aleatorios.
\end{itemize}

La robustez no elimina la necesidad de una estructura de media
correctamente especificada.
\end{frame}

\begin{frame}{21. Modelo semiparamétrico de CrossCarry}
\protect\phantomsection\label{modelo-semiparamuxe9trico-de-crosscarry}
El modelo general utilizado para medidas repetidas es

\[
g(\mu_{ijk})
=
\mathbf x_{ijk}^{\top}\boldsymbol\beta
+f(t_{ijk})
+\sum_{c=1}^{C}I_{ijc}f_c(t_{ijk}).
\]

\begin{itemize}
\tightlist
\item
  \(\mathbf x_{ijk}^{\top}\boldsymbol\beta\): tratamiento, período,
  secuencia, línea basal y covariables;
\item
  \(f(t)\): tendencia temporal común;
\item
  \(I_{ijc}\): indicador de la transición activa;
\item
  \(f_c(t)\): carry-over específico de la transición;
\item
  \(g(\cdot)\): enlace apropiado para la respuesta.
\end{itemize}

Este modelo evita imponer que el tiempo y el carry-over sigan líneas
rectas o polinomios de bajo grado.

\alert{El propósito no es hacer el modelo complejo por sí mismo, sino separar fuentes de variación que tienen interpretaciones científicas distintas.}
\end{frame}

\begin{frame}{22. Bases funcionales y splines}
\protect\phantomsection\label{bases-funcionales-y-splines}
Las funciones desconocidas se aproximan mediante bases:

\[
f(t)
=
\sum_{\ell=1}^{d}\theta_\ell\phi_\ell(t),
\qquad
f_c(t)
=
\sum_{\ell=1}^{d}\theta_{c\ell}\phi_\ell(t).
\]

Pueden emplearse:

\begin{itemize}
\tightlist
\item
  B-splines para curvas suaves y locales;
\item
  bases de Fourier para patrones periódicos;
\item
  otras bases con propiedades de aproximación adecuadas.
\end{itemize}

La base transforma la estimación funcional en la estimación de
coeficientes. El número \(d\) controla la flexibilidad:

\begin{itemize}
\tightlist
\item
  \(d\) pequeño puede producir sesgo por falta de flexibilidad;
\item
  \(d\) grande puede generar sobreajuste e inestabilidad;
\item
  la penalización permite utilizar una base razonablemente rica sin
  seguir todas las fluctuaciones de los datos.
\end{itemize}
\end{frame}

\begin{frame}{23. Dos escalas de correlación}
\protect\phantomsection\label{dos-escalas-de-correlaciuxf3n}
Con medidas repetidas existen al menos dos mecanismos de dependencia:

\begin{enumerate}
\tightlist
\item
  correlación entre tiempos dentro del mismo período;
\item
  correlación entre períodos de la misma unidad.
\end{enumerate}

Una estructura única AR(1) o intercambiable sobre las \(PL\)
observaciones puede tratar ambas escalas como si fueran equivalentes.

CrossCarry permite utilizar

\[
\mathbf R(\boldsymbol\alpha)
=
\boldsymbol\Psi\otimes
\mathbf R_1(\boldsymbol\alpha_1),
\]

con:

\begin{itemize}
\tightlist
\item
  \(\mathbf R_1\): dependencia intraperíodo;
\item
  \(\boldsymbol\Psi\): dependencia entre períodos.
\end{itemize}

La matriz resultante tiene dimensión \(PL\times PL\), pero se construye
con un número moderado de parámetros y conserva una interpretación por
bloques.
\end{frame}

\begin{frame}{24. Interpretación del producto de Kronecker}
\protect\phantomsection\label{interpretaciuxf3n-del-producto-de-kronecker}
Si

\[
\boldsymbol\Psi
=
\begin{pmatrix}
1 & \psi_{12}\\
\psi_{12} & 1
\end{pmatrix},
\]

y \(\mathbf R_1\) describe la relación entre \(L\) tiempos, entonces

\[
\boldsymbol\Psi\otimes\mathbf R_1
=
\begin{pmatrix}
\mathbf R_1 & \psi_{12}\mathbf R_1\\
\psi_{12}\mathbf R_1 & \mathbf R_1
\end{pmatrix}.
\]

Los bloques diagonales representan la correlación dentro de cada
período. Los bloques fuera de la diagonal representan la relación entre
períodos, escalada por \(\psi_{12}\).

Esta forma es útil cuando el patrón temporal intraperíodo se repite
aproximadamente en cada período, pero su intensidad entre períodos
cambia.
\end{frame}

\begin{frame}{25. Aplicación: presión arterial}
\protect\phantomsection\label{aplicaciuxf3n-presiuxf3n-arterial}
El ejemplo de presión arterial contiene:

\begin{itemize}
\tightlist
\item
  tres tratamientos: dos dosis y placebo;
\item
  seis secuencias de tres períodos;
\item
  dos individuos por secuencia;
\item
  diez mediciones dentro de cada período;
\item
  tiempos antes y después de la administración.
\end{itemize}

La respuesta temporal presenta cambios que no se describen bien con una
única línea recta. Además, no siempre es razonable asumir que un placebo
elimina el residuo de un tratamiento antihipertensivo anterior.

Los análisis mostraron:

\begin{itemize}
\tightlist
\item
  una tendencia temporal marcada;
\item
  diferencias entre algunas transiciones;
\item
  evidencia de que los carry-over no eran iguales entre tratamientos;
\item
  mejor ajuste de estructuras de correlación de Kronecker según QIC.
\end{itemize}

Este caso motivó tanto el componente funcional como la estructura de
correlación especializada.
\end{frame}

\begin{frame}{26. Aplicación: glucemia en conejas}
\protect\phantomsection\label{aplicaciuxf3n-glucemia-en-conejas}
El estudio de glucemia incluyó:

\begin{itemize}
\tightlist
\item
  dos tratamientos para control de diabetes;
\item
  secuencias ABAB y BABA;
\item
  once conejas por secuencia;
\item
  cuatro períodos;
\item
  cinco mediciones dentro de cada período.
\end{itemize}

La distribución normal no describió adecuadamente los residuos. La
comparación mediante QIC favoreció un modelo gamma con enlace
logarítmico.

El modelo permitió estimar:

\begin{itemize}
\tightlist
\item
  una disminución inicial de la glucemia seguida de recuperación;
\item
  funciones carry-over de \(A\to B\) y \(B\to A\);
\item
  ausencia de un efecto directo significativo entre tratamientos;
\item
  un efecto particular en el cuarto período.
\end{itemize}

El ejemplo ilustra que CrossCarry no se limita a respuestas gaussianas.
\end{frame}

\begin{frame}{27. Aplicación: dietas en vacas lecheras}
\protect\phantomsection\label{aplicaciuxf3n-dietas-en-vacas-lecheras}
En el estudio de producción animal, ocho vacas fueron asignadas a dos
secuencias:

\[
AB\qquad\text{y}\qquad BA.
\]

El tratamiento \(A\) correspondía a pasto y \(B\) a una mezcla que
incorporaba residuos de papel reciclado. Se realizaron mediciones
repetidas de producción y condición durante dos períodos.

El análisis con correlación de Kronecker encontró:

\begin{itemize}
\tightlist
\item
  dependencia intraperíodo importante;
\item
  una correlación positiva, aunque pequeña, entre períodos;
\item
  ausencia de un efecto directo significativo de dieta;
\item
  presencia de carry-over asociado con la dieta previa;
\item
  interacción temporal del efecto residual.
\end{itemize}

La conclusión práctica fue que el animal necesitaba tiempo para
recuperarse después del cambio de alimentación.
\end{frame}

\begin{frame}{28. Aplicación: ocupación en oficinas}
\protect\phantomsection\label{aplicaciuxf3n-ocupaciuxf3n-en-oficinas}
El experimento de ambientes de trabajo comparó cuatro diseños de oficina
mediante un diseño de Williams con 288 participantes.

Tratamientos:

\begin{itemize}
\tightlist
\item
  Nest;
\item
  Control;
\item
  Menu;
\item
  Campus.
\end{itemize}

Cada período duró dos semanas y la ocupación fue registrada mediante
sensores, generando 96 observaciones intraperíodo por participante.

Este conjunto de datos es especialmente informativo porque:

\begin{itemize}
\tightlist
\item
  contiene observaciones densas;
\item
  representa las doce transiciones;
\item
  presenta una tendencia diaria no lineal;
\item
  permite estimar carry-over complejo;
\item
  ofrece suficiente información para aplicar penalización específica por
  transición.
\end{itemize}
\end{frame}

\begin{frame}{29. Resultados del ejemplo de oficinas}
\protect\phantomsection\label{resultados-del-ejemplo-de-oficinas}
La comparación por QIC favoreció el modelo semiparamétrico penalizado
con carry-over complejo y correlación de trabajo independiente dentro de
las opciones evaluadas.

La tendencia temporal mostró:

\begin{itemize}
\tightlist
\item
  actividad creciente durante la mañana;
\item
  un descenso alrededor del almuerzo;
\item
  una segunda elevación durante la tarde.
\end{itemize}

Las transiciones no tuvieron el mismo comportamiento. Por ejemplo, el
residuo de Campus sobre Nest fue diferente del residuo de Campus sobre
Control o Menu.

Esto respalda la noción de carry-over complejo: no existe una única
curva residual de Campus que describa adecuadamente todas las
transiciones posteriores.

\alert{La forma de la curva aporta información científica adicional que desaparecería al resumir cada período mediante una sola media.}
\end{frame}

\begin{frame}{30. ¿Por qué penalizar las curvas?}
\protect\phantomsection\label{por-quuxe9-penalizar-las-curvas}
Al aumentar el número de tratamientos crece rápidamente el número de
funciones:

\[
(C+1)d
\]

coeficientes funcionales para una tendencia temporal y \(C\) carry-over,
si cada función utiliza \(d\) elementos de base.

La penalización introduce estabilidad:

\[
U_f^{\mathrm{pen}}
=U_f-\lambda\mathbf P_f\boldsymbol\theta,
\]

\[
U_c^{\mathrm{pen}}
=U_c-\lambda_c\mathbf P_c\boldsymbol\theta_c.
\]

\begin{itemize}
\tightlist
\item
  \(\lambda\) pequeño permite mayor flexibilidad;
\item
  \(\lambda\) grande fuerza una curva más suave;
\item
  un valor muy grande aproxima el efecto a una función simple o casi
  nula.
\end{itemize}

La penalización no sustituye replicación, pero evita que curvas
débilmente informadas sigan el ruido de la muestra.
\end{frame}

\begin{frame}{31. Selección de penalización y modelo}
\protect\phantomsection\label{selecciuxf3n-de-penalizaciuxf3n-y-modelo}
Los parámetros de suavidad pueden elegirse mediante:

\begin{itemize}
\tightlist
\item
  QIC en una rejilla de valores;
\item
  validación cruzada dejando una unidad completa fuera;
\item
  aproximaciones de validación cruzada generalizada;
\item
  agrupación de transiciones con comportamiento semejante para compartir
  penalización.
\end{itemize}

La comparación de modelos debe considerar conjuntamente:

\begin{itemize}
\tightlist
\item
  QIC, QICu o CIC;
\item
  convergencia;
\item
  estabilidad de estimaciones;
\item
  plausibilidad de las curvas;
\item
  condición de las matrices;
\item
  residuos;
\item
  incertidumbre de los efectos.
\end{itemize}

\alert{El menor QIC no convierte automáticamente un modelo en científicamente correcto.}
Es un criterio de comparación, no una validación completa.
\end{frame}

\begin{frame}{32. Inferencia y varianza robusta}
\protect\phantomsection\label{inferencia-y-varianza-robusta}
Para los parámetros no penalizados, la varianza robusta conserva la
forma sándwich:

\[
\widehat{\operatorname{Var}}(\widehat{\boldsymbol\beta})
=
\mathbf A^{-1}\mathbf B\mathbf A^{-1}.
\]

Para los componentes penalizados,

\[
\widehat{\operatorname{Var}}(\widehat{\boldsymbol\theta})
=
(\mathbf A+\lambda\mathbf P)^{-1}
\mathbf B
(\mathbf A+\lambda\mathbf P)^{-1}.
\]

La inferencia debe distinguir entre:

\begin{itemize}
\tightlist
\item
  efectos paramétricos, como tratamientos y períodos;
\item
  coeficientes de base;
\item
  funciones evaluadas sobre una rejilla temporal;
\item
  bandas puntuales o simultáneas;
\item
  evidencia estadística y relevancia científica.
\end{itemize}

Una curva cuyo intervalo incluye el valor nulo en casi todo el dominio
puede considerarse débil, aunque presente pequeñas oscilaciones
estimadas.
\end{frame}

\begin{frame}{33. Simulaciones: lecciones principales}
\protect\phantomsection\label{simulaciones-lecciones-principales}
Los estudios de simulación de los artículos evaluaron diseños AB/BA,
diseños de Williams, respuestas normales y Poisson, distintas
correlaciones, tamaños muestrales y números de mediciones.

Patrones recurrentes:

\begin{itemize}
\tightlist
\item
  ignorar carry-over puede deteriorar la cobertura de los intervalos;
\item
  asumir carry-over simple cuando el verdadero efecto es complejo no
  garantiza una mejora;
\item
  aumentar \(L\) mejora la reconstrucción de las funciones;
\item
  aumentar el número de unidades mejora la inferencia entre sujetos;
\item
  la correlación de Kronecker puede aumentar eficiencia cuando refleja
  la estructura real;
\item
  la penalización mejora estabilidad frente a una base no penalizada;
\item
  con muy pocas mediciones por período, el modelo funcional complejo
  puede no ser recomendable.
\end{itemize}
\end{frame}

\begin{frame}[fragile]{34. Diagnóstico del modelo}
\protect\phantomsection\label{diagnuxf3stico-del-modelo}
El diagnóstico debe examinar más que la significación de los
coeficientes.

Aspectos esenciales:

\begin{itemize}
\tightlist
\item
  residuos ordinarios, Pearson o de desviación;
\item
  patrones residuales frente al tiempo;
\item
  observaciones influyentes;
\item
  convergencia del algoritmo;
\item
  valores extremos de correlación;
\item
  sensibilidad a la dimensión de la base;
\item
  sensibilidad a la penalización;
\item
  comparación de estructuras de trabajo.
\end{itemize}

Código conceptual:

\begin{Shaded}
\begin{Highlighting}[]
\CommentTok{\# Objetos y nombres exactos dependen de la función ajustada}
\FunctionTok{summary}\NormalTok{(modelo)}
\FunctionTok{plot}\NormalTok{(modelo)}
\NormalTok{modelo}\SpecialCharTok{$}\NormalTok{QIC}
\end{Highlighting}
\end{Shaded}

Los residuos no validan por sí solos el diseño ni demuestran ausencia de
carry-over. Sirven para detectar incompatibilidades importantes entre
datos y modelo.
\end{frame}

\begin{frame}[fragile]{35. Arquitectura del paquete CrossCarry}
\protect\phantomsection\label{arquitectura-del-paquete-crosscarry}
El paquete organiza el análisis en cuatro funciones principales:

\begin{itemize}
\tightlist
\item
  \texttt{createCarry()}: codificación de carry-over de primer orden;
\item
  \texttt{CrossGEE()}: GEE con estructuras convencionales;
\item
  \texttt{CrossGEEKron()}: GEE con correlación por producto de
  Kronecker;
\item
  \texttt{CrossGEESP()}: GEE semiparamétrica para tiempo y carry-over.
\end{itemize}

Flujo recomendado:

\begin{enumerate}
\tightlist
\item
  organizar y verificar los datos;
\item
  construir las variables carry-over;
\item
  ajustar modelos coherentes con el diseño;
\item
  comparar estructura de media y correlación;
\item
  revisar diagnóstico e incertidumbre;
\item
  interpretar en la escala del enlace.
\end{enumerate}

No todas las funciones admiten las mismas opciones. La versión instalada
y la documentación del paquete deben revisarse antes de ejecutar un
análisis reproducible.
\end{frame}

\begin{frame}[fragile]{36. Preparación del ejemplo Arterial}
\protect\phantomsection\label{preparaciuxf3n-del-ejemplo-arterial}
\begin{Shaded}
\begin{Highlighting}[]
\FunctionTok{library}\NormalTok{(CrossCarry)}
\FunctionTok{data}\NormalTok{(}\StringTok{"Arterial"}\NormalTok{)}

\FunctionTok{str}\NormalTok{(Arterial)}
\FunctionTok{head}\NormalTok{(Arterial)}
\end{Highlighting}
\end{Shaded}

La base contiene las variables necesarias para el análisis:

\begin{itemize}
\tightlist
\item
  \texttt{Pressure}: presión arterial;
\item
  \texttt{Subject}: identificador individual;
\item
  \texttt{Treatment}: tratamiento del período;
\item
  \texttt{Period}: período del diseño;
\item
  \texttt{Time}: momento de medición dentro del período.
\end{itemize}

Antes de ajustar:

\begin{Shaded}
\begin{Highlighting}[]
\FunctionTok{with}\NormalTok{(Arterial, }\FunctionTok{table}\NormalTok{(Treatment, Period))}
\FunctionTok{with}\NormalTok{(Arterial, }\FunctionTok{table}\NormalTok{(Subject, Period))}
\FunctionTok{summary}\NormalTok{(Arterial}\SpecialCharTok{$}\NormalTok{Time)}
\end{Highlighting}
\end{Shaded}

Estas comprobaciones permiten detectar ausencias, duplicados y
desequilibrios que podrían afectar la construcción del carry-over o la
correlación.
\end{frame}

\begin{frame}[fragile]{37. Construcción de indicadores carry-over}
\protect\phantomsection\label{construcciuxf3n-de-indicadores-carry-over}
La interfaz documentada en el manuscrito del paquete utiliza:

\begin{Shaded}
\begin{Highlighting}[]
\NormalTok{carrydata }\OtherTok{\textless{}{-}} \FunctionTok{createCarry}\NormalTok{(}
  \AttributeTok{data =}\NormalTok{ Arterial,}
  \AttributeTok{treatment =} \StringTok{"Treatment"}\NormalTok{,}
  \AttributeTok{period =} \StringTok{"Period"}\NormalTok{,}
  \AttributeTok{id =} \StringTok{"Subject"}\NormalTok{,}
  \AttributeTok{carrySimple =} \ConstantTok{TRUE}
\NormalTok{)}

\NormalTok{datos }\OtherTok{\textless{}{-}}\NormalTok{ carrydata}\SpecialCharTok{$}\NormalTok{data}
\NormalTok{carry }\OtherTok{\textless{}{-}}\NormalTok{ carrydata}\SpecialCharTok{$}\NormalTok{carryover}

\FunctionTok{head}\NormalTok{(datos)}
\FunctionTok{head}\NormalTok{(carry)}
\end{Highlighting}
\end{Shaded}

Con carry-over simple, una columna como \texttt{Carry\_B} vale 1 cuando
el tratamiento inmediatamente anterior fue \(B\) y 0 en otro caso.

Para carry-over complejo, la codificación debe representar transiciones
como \(A\to B\) o \(C\to A\). Antes de modelarlas debe comprobarse que
cada columna contenga observaciones activas suficientes.
\end{frame}

\begin{frame}[fragile]{38. Ajuste con correlación de Kronecker}
\protect\phantomsection\label{ajuste-con-correlaciuxf3n-de-kronecker}
\begin{Shaded}
\begin{Highlighting}[]
\NormalTok{modelo\_kron }\OtherTok{\textless{}{-}} \FunctionTok{CrossGEEKron}\NormalTok{(}
  \AttributeTok{response =} \StringTok{"Pressure"}\NormalTok{,}
  \AttributeTok{treatment =} \StringTok{"Treatment"}\NormalTok{,}
  \AttributeTok{period =} \StringTok{"Period"}\NormalTok{,}
  \AttributeTok{id =} \StringTok{"Subject"}\NormalTok{,}
  \AttributeTok{time =} \StringTok{"Time"}\NormalTok{,}
  \AttributeTok{carry =}\NormalTok{ carry,}
  \AttributeTok{data =}\NormalTok{ datos,}
  \AttributeTok{correlation =} \StringTok{"AR{-}M"}\NormalTok{,}
  \AttributeTok{Mv =} \DecValTok{1}
\NormalTok{)}

\FunctionTok{summary}\NormalTok{(modelo\_kron}\SpecialCharTok{$}\NormalTok{model)}
\NormalTok{modelo\_kron}\SpecialCharTok{$}\NormalTok{QIC}
\end{Highlighting}
\end{Shaded}

El resumen debe interpretarse mediante los errores estándar robustos.
Los coeficientes de tratamiento comparan cada nivel con la categoría de
referencia, condicionados por período, tiempo y carry-over incluidos en
el predictor.

La matriz de Kronecker permite recuperar:

\begin{itemize}
\tightlist
\item
  el patrón dentro del período;
\item
  la relación entre períodos;
\item
  la matriz completa de trabajo.
\end{itemize}
\end{frame}

\begin{frame}[fragile]{39. Selección, interpretación y reporte}
\protect\phantomsection\label{selecciuxf3n-interpretaciuxf3n-y-reporte}
Una comparación reproducible puede organizarse así:

\begin{Shaded}
\begin{Highlighting}[]
\NormalTok{comparacion }\OtherTok{\textless{}{-}} \FunctionTok{data.frame}\NormalTok{(}
  \AttributeTok{modelo =} \FunctionTok{c}\NormalTok{(}
    \StringTok{"Independencia"}\NormalTok{, }\StringTok{"AR(1)"}\NormalTok{,}
    \StringTok{"Kronecker"}\NormalTok{, }\StringTok{"Semiparamétrico"}
\NormalTok{  ),}
  \AttributeTok{QIC =} \FunctionTok{c}\NormalTok{(qic\_ind, qic\_ar1, qic\_kron, qic\_sp)}
\NormalTok{)}

\NormalTok{comparacion[}\FunctionTok{order}\NormalTok{(comparacion}\SpecialCharTok{$}\NormalTok{QIC), ]}
\end{Highlighting}
\end{Shaded}

El informe final debe incluir:

\begin{itemize}
\tightlist
\item
  diseño y secuencias;
\item
  duración de períodos y lavado;
\item
  número de unidades y pérdidas;
\item
  tiempos de medición;
\item
  familia y enlace;
\item
  codificación del carry-over;
\item
  estructura de correlación;
\item
  dimensión de base y penalización;
\item
  estimaciones robustas y bandas;
\item
  diagnóstico y análisis de sensibilidad.
\end{itemize}
\end{frame}

\begin{frame}{40. Conclusiones y preguntas}
\protect\phantomsection\label{conclusiones-y-preguntas}
CrossCarry integra una trayectoria metodológica para diseños crossover
que incluye:

\begin{itemize}
\tightlist
\item
  respuestas gaussianas y no gaussianas;
\item
  una o múltiples mediciones por período;
\item
  carry-over simple y complejo;
\item
  tendencias temporales no lineales;
\item
  correlación dentro y entre períodos;
\item
  estimación semiparamétrica y penalizada;
\item
  criterios de comparación y herramientas de interpretación.
\end{itemize}

Recomendaciones finales:

\begin{enumerate}
\tightlist
\item
  diseñar pensando en las transiciones que se desean estimar;
\item
  justificar científicamente el período de lavado;
\item
  no confundir muchos registros con suficiente replicación;
\item
  comenzar por modelos interpretables y aumentar complejidad solo cuando
  los datos la sostengan;
\item
  interpretar efectos en su escala científica, no solo mediante valores
  \(p\).
\end{enumerate}

\vfill

\centering

\Large\textbf{CrossCarry convierte la memoria del tratamiento en una parte explícita, estimable e interpretable del análisis.}

\vspace{0.4cm}

\textbf{Gracias por su atención. ¿Preguntas?}
\end{frame}




\end{document}
